\section{Rappels de mécanique des matériaux}\label{sec:meca}

\subsection{Petites déformations}

Soit $\mathbf{u}(\mathbf{x})$ le champ de déplacement. Sous l'hypothèse des petites perturbations ($\|\nabla\mathbf{u}\|\ll1$), la déformation est mesurée par la partie symétrique du gradient :
\begin{equation}
\eps = \tfrac12\big(\nabla\mathbf{u}+\nabla\mathbf{u}^T\big).
\end{equation}
$\tr\eps$ mesure la variation relative de volume et $\dev\eps$ le changement de forme. En élastoplasticité on décompose la déformation de façon \emph{additive} :
\begin{equation}
\eps = \eps^e + \epsp ,
\end{equation}
où $\eps^e$ est réversible (élastique) et $\epsp$ permanente (plastique).

\subsection{Élasticité linéaire isotrope}

La loi de Hooke isotrope s'écrit avec les coefficients de Lamé $(\lambda,\mu)$ :
\begin{equation}
\sig = \lambda\tr(\eps^e)\,\I + 2\mu\,\eps^e = \C:\eps^e,\qquad \C = 3K\,\Id_{\mathrm{vol}} + 2\mu\,\Id_{\mathrm{dev}},
\label{eq:hooke}
\end{equation}
où $K=\lambda+\frac23\mu$ est le module de compressibilité. La forme projetée est la plus utile ici : elle découple le volume et la forme, $\tr\sig = 3K\tr\eps^e$ et $\dev\sig = 2\mu\dev\eps^e$. En fonction du module d'Young $E$ et du coefficient de Poisson $\nu$ :
\[
\mu = \frac{E}{2(1+\nu)},\qquad K = \frac{E}{3(1-2\nu)},\qquad \lambda = \frac{E\nu}{(1+\nu)(1-2\nu)}.
\]

\subsection{Critère de von Mises}

\begin{rappel}{pourquoi J2}
Dans les métaux, la plasticité vient du glissement des dislocations, qui est piloté par le cisaillement. Une pression hydrostatique ne fait pas plastifier le matériau, et l'écoulement plastique se fait sans variation de volume. Le critère ne doit donc dépendre que du déviateur $\s=\dev\sig$. Le plus simple de ses invariants est $J_2 = \frac12\s:\s$, d'où le nom « plasticité J2 ».
\end{rappel}

La contrainte équivalente de von Mises est normalisée pour valoir $|\sigma|$ en traction simple :
\begin{equation}
\sigma_{eq}(\sig) = \sqrt{3J_2} = \sqrt{\tfrac32\,\s:\s}.
\end{equation}
\emph{Vérification :} pour $\sig=\sigma\,\mathbf{e}_1\otimes\mathbf{e}_1$, on a $\s = \frac{\sigma}{3}\mathrm{diag}(2,-1,-1)$, donc $\s:\s=\frac{6\sigma^2}{9}$ et $\sigma_{eq}=|\sigma|$.

Le domaine d'élasticité est $\{\sig : f(\sig,\ldots)\le0\}$, avec la fonction de charge
\begin{equation}
f(\sig,\X,p) = \sigma_{eq}(\sig-\X) - R(p),
\label{eq:f}
\end{equation}
où $\X$ (tenseur de rappel, déviatorique) est la position du centre du domaine et $R(p)$ son rayon (au sens de $\sigma_{eq}$).

\subsection{Écoulement plastique et déformation cumulée}

L'écoulement est \emph{associé} (normal au domaine) :
\begin{equation}
\dot{\epsp} = \dot p\,\N,\qquad \N = \frac{\partial f}{\partial\sig} = \frac32\,\frac{\bs\xi}{\sigma_{eq}(\bs\xi)},\qquad \bs\xi = \s-\X .
\label{eq:flow}
\end{equation}
Comme $\N:\N=\frac94\frac{\bs\xi:\bs\xi}{\frac32\bs\xi:\bs\xi}=\frac32$, le multiplicateur $\dot p$ est exactement la \emph{déformation plastique cumulée équivalente} : $\dot p = \sqrt{\tfrac23\dot{\epsp}:\dot{\epsp}}$. Le flux $\N$ est déviatorique, donc $\tr\dot{\epsp}=0$ : l'écoulement est isochore.

\subsection{Écrouissages}

\paragraph{Isotrope.} Le domaine grossit avec $p$. On utilise
\begin{equation}
R(p) = \sigma_y + H\,p + Q\big(1-e^{-bp}\big),\qquad R'(p) = H + Qb\,e^{-bp}.
\end{equation}
Pour $Q=0$ l'écrouissage est linéaire (cas du sujet). Pour $Q>0$ c'est la loi de Voce~\cite{voce1948}, qui sature vers $\sigma_y+Q$ avec une vitesse $b$. C'est le comportement typique des aciers recuits.

\paragraph{Cinématique (Prager~\cite{prager1956}).} Le domaine se translate, $\dot{\X} = \frac23\,C\,\dot{\epsp}$. Après une traction plastique, $\X$ pointe dans le sens de la traction. La plastification en compression commence alors \emph{avant} $-\sigma_y$ : c'est l'\emph{effet Bauschinger}, indispensable pour la fatigue et les chargements cycliques. On retrouvera cet effet sur les boucles de la figure~\ref{fig:hyst}.

\subsection{Conditions de Kuhn--Tucker et cohérence}

Le multiplicateur est déterminé par les conditions de complémentarité
\begin{equation}
f\le0,\qquad \dot p\ge0,\qquad \dot p\,f=0,
\end{equation}
qui disent qu'il n'y a pas d'écoulement à l'intérieur du domaine, et qu'en cas d'écoulement l'état reste sur la frontière. Pendant l'écoulement, la condition de \emph{cohérence} $\dot f=0$ détermine $\dot p$.

\subsection{Réponse uniaxiale analytique}\label{sec:uniax}

En traction simple ($\sigma_{11}=\sigma$, autres composantes nulles) avec écrouissage linéaire total $h=H+C$, on a $\X = \frac{2}{3}C\epsp$ et $\epsp = p\,\mathrm{diag}(1,-\frac12,-\frac12)$. La condition $f=0$ donne $\sigma - Cp = \sigma_y + Hp$, d'où $\sigma=\sigma_y+h\,p$. Avec $\varepsilon = \sigma/E + p$, on obtient
\begin{equation}
\sigma = \sigma_y + E_t\,(\varepsilon-\varepsilon_y),\qquad E_t = \frac{E\,h}{E+h},\qquad \varepsilon_y=\frac{\sigma_y}{E}.
\label{eq:Et}
\end{equation}
C'est la solution de référence du test \code{test\_uniaxial\_bilinear\_analytic}.
